\documentclass[../../main.tex]{subfiles}
\begin{document}
\section{群作用的无穷小生成元}\label{section:ch20-infinitesimal-generator}

在第 9 章中，我们证明了流形上一个完全向量场生成 $\mathbb{R}$ 在该流形上的一个作用。本节中，利用 Frobenius 定理与指数映射的性质，我们说明如何把这一概念推广到更高维群的作用。

首先，我们需要明确所谓 Lie 群作用的"无穷小生成元"是什么意思。由于后面将说明的原因，本节主要讨论右作用。之后我们会说明在左作用情形理论应如何修改（见定理\cref{theorem:ch20-左作用基本定理}）。由于 $\mathbb{R}$ 交换，整体流既可视为左作用也可视为右作用，故本节中一切对整体流无需修改即适用。

设给定 Lie 群 $G$ 在光滑流形 $M$ 上的一个光滑右作用，依上下文记为 $\sigma\colon M\times G\to M$ 或 $(p,g)\mapsto pg$。每个元素 $X\in\operatorname{Lie}(G)$ 确定 $M$ 上的一个光滑整体流
\begin{align*}
(t,p)\mapsto p\cdot\exp tX.
\end{align*}
（它是流，因为 $\exp 0X=e$ 且 $(\exp sX)(\exp tX)=\exp(s+t)X$。）令 $\widetilde X\in\mathfrak X(M)$ 为该流的无穷小生成元，于是对每个 $p\in M$，
\begin{align}
\widetilde X_p=\frac{\mathrm{d}}{\mathrm{d}t}\bigg|_{t=0}p\cdot\exp tX.
\end{align}
于是我们得到映射 $\widetilde{(\,)}\colon\operatorname{Lie}(G)\to\mathfrak X(M)$，由 $\widetilde{(X)}=\widetilde X$ 定义。

$\widetilde X$ 还有一个用轨道映射 $\sigma_p\colon G\to M$（$\sigma_p(g)=p\cdot g$）给出的有用等价刻画。由于 $\gamma(t)=\exp tX$ 是 $G$ 中初速度 $\gamma'(0)=X_e$ 的光滑曲线，由推论\cref{corollary:ch03-25} 得对每个 $p\in M$，
\begin{align}
\mathrm{d}(\sigma_p)_e(X_e)=\sigma_p\circ\gamma'(0)=\frac{\mathrm{d}}{\mathrm{d}t}\bigg|_{t=0}p\cdot\exp tX=\widetilde X_p.
\end{align}

\begin{lemma}\label{lemma:ch20-轨道映射相关}
设 $G$ 是一个 Lie 群，$\sigma$ 是 $G$ 在光滑流形 $M$ 上的一个光滑右作用。对任意 $X\in\operatorname{Lie}(G)$ 与 $p\in M$，向量场 $X$ 与 $\widetilde{(X)}$ 是 $\sigma_p$-相关的。
\end{lemma}

\begin{proposition}\label{proposition:ch20-生成元是同态}
映射 $\widetilde{(\,)}\colon\operatorname{Lie}(G)\to\mathfrak X(M)$ 是一个 Lie 代数同态。
\end{proposition}

\begin{proof}
由\eqref{lemma:ch20-轨道映射相关}，$\widetilde X_p$ 线性依赖于 $X$，故 $\widetilde{(\,)}$ 是线性映射。给定 $p\in M$，引理\cref{lemma:ch20-轨道映射相关} 连同 Lie 括号的自然性蕴含 $[X,Y]$ 与 $[\widetilde X,\widetilde Y]$ 是 $\sigma_p$-相关的。这特别意味着
\begin{align*}
[\widetilde X,\widetilde Y]_p=\mathrm{d}(\sigma_p)_e\bigl([X,Y]_e\bigr)=\widetilde{[X,Y]}_p.
\end{align*}
由于 $M$ 的每一点都在某个轨道映射的像中，即得 $[\widetilde{(X)},\widetilde{(Y)}]=\widetilde{([X,Y])}$。
\end{proof}

上述定义的 Lie 代数同态 $\widetilde{(\,)}\colon\operatorname{Lie}(G)\to\mathfrak X(M)$ 称为 $\sigma$ 的\emph{无穷小生成元}。更一般地，若 $\mathfrak g$ 是任意有限维 Lie 代数，任意 Lie 代数同态 $\widetilde{(\,)}\colon\mathfrak g\to\mathfrak X(M)$ 称为 $M$ 上的一个（右）$\mathfrak g$-作用。若对每个 $X\in\mathfrak g$，向量场 $\widetilde{(X)}$ 都是完全的，则称该 $\mathfrak g$-作用是\emph{完全的}。

正如每个完全向量场生成一个 $\mathbb{R}$-作用，下一定理表明，至少对单连通群而言，每个完全 Lie 代数作用都生成一个 Lie 群作用。

\begin{theorem}{[Lie 代数作用基本定理]}\label{theorem:ch20-Lie代数作用基本定理}
设 $M$ 是一个光滑流形，$G$ 是一个单连通 Lie 群，$\mathfrak g=\operatorname{Lie}(G)$。设 $\widetilde{(\,)}\colon\mathfrak g\to\mathfrak X(M)$ 是 $M$ 上的一个完全 $\mathfrak g$-作用。则存在 $M$ 上唯一的光滑右 $G$-作用，其无穷小生成元为 $\widetilde{(\,)}$。
\end{theorem}

\begin{proof}
我们先在 $G\times M$ 上定义一个分布 $D$；将证明 $D$ 对合，然后每一叶将成为某个轨道映射 $\sigma_p\colon G\to M$ 的图像。为简洁起见，给定 $X\in\mathfrak g$，用记号 $\widetilde X$ 表示 $\widetilde{(X)}\in\mathfrak X(M)$。

对每个 $X\in\mathfrak g$，在 $G\times M$ 上定义光滑向量场
\begin{align*}
\zeta_{(X)}(g,p)=\bigl(X_g,\widetilde X_p\bigr).
\end{align*}
令 $D$ 为由所有这些 $\zeta_{(X)}$ 张成的分布。

首先证 $D$ 对合。对 $X,Y\in\mathfrak g$，直接计算表明 $[\zeta_{(X)},\zeta_{(Y)}]=\zeta_{([X,Y])}$：其 $G$ 分量为 $[X,Y]_g$（左不变向量场的括号仍是左不变的），其 $M$ 分量由命题\cref{proposition:ch20-生成元是同态} 为 $[\widetilde X,\widetilde Y]_p=\widetilde{[X,Y]}_p$。因此 $D$ 对合。

令 $\mathscr F$ 为 $D$ 确定的叶状结构。注意对每个 $g_0\in G$，左平移 $L_{g_0}$ 给出 $G\times M$ 的一个自同胚 $(g,p)\mapsto(g_0g,p)$。直接计算表明
\begin{align*}
\mathrm{d}(L_{g_0})_{(g_0^{-1}g,p)}\bigl(\zeta_{(X)}(g_0^{-1}g,p)\bigr)=\zeta_{(X)}(g,p),
\end{align*}
故 $D$ 在 $L_{g_0}$ 下不变，从而 $L_{g_0}$ 把 $\mathscr F$ 的叶映为 $\mathscr F$ 的叶（命题\cref{proposition:ch19-不变分布与叶状}）。

令 $\pi_G\colon G\times M\to G$、$\pi_M\colon G\times M\to M$ 为投影。任取 $p\in M$，令 $S_p$ 为含 $(e,p)$ 的叶。我们将证 $\pi_G|_{S_p}\colon S_p\to G$ 是光滑覆叠映射。首先，在每点 $(g,q)\in S_p$，$\mathrm{d}(\pi_G)(\zeta_{(X)})=X_g$ 对所有 $X\in\mathfrak g$ 成立，故 $\pi_G|_{S_p}$ 是光滑淹没，由维数原因它是局部微分同胚。

为证它是覆叠映射，取 $e$ 的连通邻域 $U\subseteq G$ 足够小，使 $G$ 的指数映射是从 $\mathfrak g$ 中 $0$ 的某个邻域 $V$ 到 $U$ 的微分同胚；对任意 $g\in G$，考虑 $gU=\{gh\colon h\in U\}$。我们将通过构造局部截面证明 $gU$ 被均匀覆盖。对每个使 $(g,q)$ 在纤维 $(\pi_G|_{S_p})^{-1}(g)$ 中的 $q\in M$，定义映射 $\sigma_q\colon gU\to G\times M$，
\begin{align*}
\sigma_q(g\exp X)=\bigl(g\exp X,(\widetilde X)_{(1,q)}\bigr),
\end{align*}
其中 $(\widetilde X)_{(t,q)}$ 是 $\widetilde X$ 的流。这是良定义的，因为 $U$ 中 $g\exp X$ 的表示唯一，且 $(g,q)\in S_p$ 保证 $(g\exp X,(\widetilde X)_{(1,q)})$ 仍在 $S_p$ 中。这些局部截面彼此不交、覆盖整个纤维、并都映到 $gU$。故 $\pi_G|_{S_p}$ 是覆叠映射。

由于 $G$ 单连通，连通覆叠空间 $S_p$ 到 $G$ 的覆叠映射必为微分同胚。于是存在唯一的光滑映射 $\sigma_p\colon G\to M$，使 $S_p=\{(g,\sigma_p(g))\colon g\in G\}$。

定义右作用如下：对 $p\in M$、$g\in G$，令 $p\cdot g=\sigma_p(g)$。需证这与基点 $p$ 的选取一致，即若 $(g,q)\in S_p$ 则 $\sigma_q(g_0)=\sigma_p(gg_0)$。等价地，设 $S_{(e,p)}=S_{(g,q)}$ 且 $S_{(e,q)}=S_{(g_0,r)}$，需证 $S_{(e,p)}=S_{(gg_0,r)}$。这由 $L_g$-不变性得到：
\begin{align*}
S_{(e,p)}=S_{(g,q)}=L_g\bigl(S_{(e,q)}\bigr)=L_g\bigl(S_{(g_0,r)}\bigr)=S_{(gg_0,r)}.
\end{align*}

还需证作用光滑、$\widetilde{(\,)}$ 是其无穷小生成元、且它是唯一的此类作用。对单位元附近的 $g=\exp X$，上述讨论表明作用可表为
\begin{align}
p\cdot g=\sigma_p(\exp X)=\pi_M\circ\sigma_p(\exp X)=(\widetilde X)_{(1,p)}.
\end{align}
一个与我们证明指数映射光滑所用者类似的论证（在 $M\times\mathfrak g$ 上取 $(p,X)\mapsto(\widetilde X_p,0)$）表明这光滑依赖于 $X$ 与 $p$，因而光滑依赖于 $g$ 与 $p$。由于单位元的任意邻域生成 $G$（命题\cref{proposition:ch07-14}），$G$ 的每个元素都可表为形如 $\exp X$（$X\in V$）的元素的有限乘积，故 $(p,g)\mapsto p\cdot g$ 可写成光滑映射的有限复合。作用的无穷小生成元为 $\widetilde{(\,)}$ 是上式的直接推论。唯一性留作习题。
\end{proof}

\begin{exercise}\label{exercise:ch20-17}
证明上述证明中构造的作用是以 $\widetilde{(\,)}$ 为无穷小生成元的唯一作用。
\end{exercise}

\subsection{左作用}

左作用的情形类似，但略有转折。设 $G$ 是 Lie 群，$M$ 是光滑流形。若 $\sigma\colon G\times M\to M$ 是 $G$ 在 $M$ 上的光滑左作用，定义 $\sigma$ 的无穷小生成元为映射 $\widetilde{(\,)}\colon\operatorname{Lie}(G)\to\mathfrak X(M)$，由 $\widetilde{(X)}=\widetilde X$ 给出，其中
\begin{align*}
\widetilde X_p=\frac{\mathrm{d}}{\mathrm{d}t}\bigg|_{t=0}\sigma(\exp tX,p).
\end{align*}

\begin{theorem}{[左 Lie 代数作用基本定理]}\label{theorem:ch20-左作用基本定理}
设 $M$ 是光滑流形，$G$ 是单连通 Lie 群，$\mathfrak g=\operatorname{Lie}(G)$。设 $\widehat{(\,)}\colon\mathfrak g\to\mathfrak X(M)$ 是一个 Lie 代数反同态（即 $[\widehat{(X)},\widehat{(Y)}]=-\widehat{([X,Y])}$），且对每个 $X$，$\widehat{(X)}$ 完全。则存在 $M$ 上唯一的光滑左 $G$-作用，其无穷小生成元为 $\widehat{(\,)}$。
\end{theorem}

由于此定理，对有限维 Lie 代数 $\mathfrak g$ 与光滑流形 $M$，$M$ 上的左 $\mathfrak g$-作用定义为从 $\mathfrak g$ 到 $\mathfrak X(M)$ 的一个反同态。
\end{document}
